\begingroup\let\clearpageforsection\relax\exercisesheader{}\endgroup

% 25

\eoce{\qt{Efek samping Avandia\label{randomization_avandia}} Rosiglitazon adalah
bahan aktif dalam obat diabetes tipe~2 Avandia yang kontroversial dan telah
dikaitkan dengan peningkatan risiko masalah kardiovaskular serius, seperti
stroke, gagal jantung, dan kematian. Salah satu pengobatan alternatif yang umum
adalah pioglitazon, yaitu bahan aktif dalam obat diabetes bernama Actos. Dalam
sebuah studi observasional retrospektif berskala nasional terhadap 227,571
penerima manfaat Medicare berusia 65 tahun atau lebih, ditemukan bahwa 2,593
dari 67,593 pasien yang menggunakan rosiglitazon dan 5,386 dari 159,978 pasien
yang menggunakan pioglitazon mengalami masalah kardiovaskular serius. Data ini
dirangkum dalam tabel kontingensi berikut. \footfullcite{Graham:2010}
\begin{center}
\begin{tabular}{ll  cc c} 
                                &   & \multicolumn{2}{c}{\textit{Masalah kardiovaskular}} \\
\cline{3-4} 
                                    &               & Ya    & Tidak     & Total \\
\cline{2-5}
\multirow{2}{*}{\textit{Perlakuan}} & Rosiglitazon  & 2,593 & 65,000    & 67,593 \\
                                    & Pioglitazon   & 5,386 & 154,592   & 159,978 \\
\cline{2-5}
                                    & Total         & 7,979 & 219,592   & 227,571
\end{tabular}
\end{center}
\begin{parts}
\item Tentukan apakah setiap pernyataan berikut benar atau salah. Jika salah,
jelaskan alasannya. \textit{Berhati-hatilah:} Penalarannya mungkin salah walaupun
kesimpulan pernyataannya benar. Dalam kasus seperti itu, pernyataan harus
dianggap salah.
\begin{subparts}
\item Karena lebih banyak pasien pengguna pioglitazon mengalami masalah
kardiovaskular (5,386 dibandingkan dengan 2,593), kita dapat menyimpulkan bahwa
tingkat masalah kardiovaskular pada orang yang menerima perlakuan pioglitazon
lebih tinggi.
\item Data mengisyaratkan bahwa pasien diabetes yang menggunakan rosiglitazon
lebih mungkin mengalami masalah kardiovaskular karena tingkat kejadiannya
adalah (2,593 / 67,593 = 0.038) 3.8\% pada pasien yang menerima perlakuan ini,
sedangkan pada pasien pengguna pioglitazon hanya
(5,386 / 159,978 = 0.034) 3.4\%.
\item Fakta bahwa tingkat kejadian lebih tinggi pada kelompok rosiglitazon
membuktikan bahwa rosiglitazon menyebabkan masalah kardiovaskular serius.
\item Berdasarkan informasi yang diberikan sejauh ini, kita tidak dapat
menentukan apakah perbedaan tingkat kejadian disebabkan oleh hubungan antara
kedua variabel atau oleh faktor kebetulan.
\end{subparts}
\item Berapa proporsi seluruh pasien yang mengalami masalah kardiovaskular?
\item Jika jenis perlakuan dan kejadian masalah kardiovaskular saling
independen, kira-kira berapa banyak pasien dalam kelompok rosiglitazon yang
diperkirakan mengalami masalah kardiovaskular?
\item Kita dapat menyelidiki hubungan antara hasil dan perlakuan dalam penelitian
ini dengan teknik pengacakan. Dalam praktiknya, simulasi yang diperlukan untuk
pengacakan akan dijalankan menggunakan perangkat lunak statistika, tetapi
misalkan kita benar-benar menyimulasikannya dengan kartu indeks. Untuk
menyimulasikan model independensi, yang menyatakan bahwa hasil tidak bergantung
pada perlakuan, pada setiap kartu kita tuliskan apakah seorang pasien mengalami
masalah kardiovaskular atau tidak. Semua kartu kemudian dikocok dan dibagikan
ke dalam dua kelompok berukuran 67,593 dan 159,978. Kita mengulangi simulasi ini
1,000 kali dan setiap kali mencatat jumlah orang dalam kelompok rosiglitazon
yang mengalami masalah kardiovaskular. Gunakan histogram frekuensi relatif
dari jumlah tersebut untuk menjawab (i)--(iii).
\end{parts}
\begin{minipage}[c]{0.5\textwidth}
\begin{subparts}
\item Pernyataan apa saja yang sedang diuji?
\item Dibandingkan dengan jumlah yang dihitung pada bagian~(c), manakah yang
akan lebih mendukung hipotesis alternatif: \textit{lebih banyak} atau
\textit{lebih sedikit} pasien dengan masalah kardiovaskular dalam kelompok
rosiglitazon?
\item Apa yang ditunjukkan oleh hasil simulasi mengenai hubungan antara
penggunaan rosiglitazon dan kejadian masalah kardiovaskular pada pasien diabetes?
\end{subparts}
\end{minipage}
\begin{minipage}[c]{0.5\textwidth}
\Figures[Ditampilkan histogram "Kejadian kardiovaskular rosiglitazon hasil simulasi", dengan nilai yang membentang dari 2250 hingga 2450. Dari sisi kiri, histogram dimulai dengan interval-interval kelas yang rendah sampai sekitar 2280. Tinggi kelas kemudian naik bertahap, lalu meningkat tajam mulai 2320 hingga mencapai puncak sekitar 2360. Tinggi kelas turun tajam sekitar 2380 menjadi sekitar setengah tinggi puncak, lalu menurun bertahap sampai 2460 dan menjadi nol setelah titik tersebut.]{}{eoce/randomization_avandia}{randomization_avandia} \\
\end{minipage}
}{}


% 26

\eoce{\qt{Transplantasi jantung\label{randomization_heart_transplants}} Studi
Transplantasi Jantung Universitas Stanford dilakukan untuk menentukan apakah
program transplantasi jantung eksperimental memperpanjang masa hidup. Setiap
pasien yang mengikuti program ditetapkan sebagai calon resmi penerima
transplantasi jantung. Artinya, pasien tersebut sakit parah dan kemungkinan
besar akan memperoleh manfaat dari jantung baru. Sebagian pasien menerima
transplantasi dan sebagian lainnya tidak. Variabel \texttt{transplant}
menunjukkan kelompok pasien; pasien dalam kelompok perlakuan menerima
transplantasi, sedangkan pasien dalam kelompok kontrol tidak. Dari 34 pasien
dalam kelompok kontrol, 30 meninggal. Dari 69 orang dalam kelompok perlakuan,
45 meninggal. Variabel lain bernama \texttt{survived} digunakan untuk menunjukkan
apakah pasien masih hidup pada akhir penelitian. \footfullcite{Turnbull+Brown+Hu:1974}
\begin{center}
\Figures[Diagram mosaik untuk variabel "kelompok eksperimen" (pemisahan primer) dan "kelangsungan hidup". Persegi panjang tinggi pertama untuk kelompok eksperimen "kontrol" lebarnya sekitar setengah persegi panjang tinggi kedua untuk "perlakuan". Pada pemisahan sekunder kelompok kontrol, hasil "hidup" mencakup sekitar 10\% tinggi persegi panjang dan "meninggal" sekitar 90\%. Pada pemisahan sekunder kelompok perlakuan, hasil "hidup" mencakup sekitar 35\% tinggi persegi panjang dan "meninggal" sekitar 65\%.]{0.48}{eoce/randomization_heart_transplants}{randomization_heart_transplants_mosaic}
\Figures[Ditampilkan diagram kotak berdampingan untuk variabel "Masa hidup (hari)", dengan dua diagram kotak berlabel "kontrol" dan "perlakuan". Sumbu masa hidup membentang dari 0 hingga sekitar 1800. Kotak kelompok kontrol membentang dari sekitar 0 hingga 50 dengan garis median sekitar 20, sedangkan kumisnya memanjang ke bawah hingga 0 dan ke atas hingga sekitar 125. Lima pengamatan berada di atas kumis atas, masing-masing sekitar 150, 250, 300, 400, dan 1400. Kotak kelompok perlakuan membentang dari sekitar 100 hingga 650 dengan garis median sekitar 250, sedangkan kumisnya memanjang ke bawah hingga 0 dan ke atas hingga sekitar 1400. Beberapa titik berada di atas kumis atas, masing-masing sekitar 1550, 1575, dan 1800.]{0.48}{eoce/randomization_heart_transplants}{randomization_heart_transplants_box}
\end{center}
\begin{parts}
\item Berdasarkan diagram mosaik, apakah kelangsungan hidup independen dari
status penerimaan transplantasi? Jelaskan penalaran Anda.
\item Apa yang ditunjukkan oleh diagram kotak di atas mengenai efikasi
(keefektifan) perlakuan transplantasi jantung?
\item Berapa proporsi pasien yang meninggal dalam kelompok perlakuan, dan berapa
proporsinya dalam kelompok kontrol?
\item Salah satu pendekatan untuk menyelidiki apakah perlakuan tersebut efektif
adalah dengan menggunakan teknik pengacakan.
\begin{subparts}
\item Pernyataan apa saja yang sedang diuji?
\item Paragraf berikut menjelaskan penyiapan pendekatan tersebut jika kita
melakukannya tanpa perangkat lunak statistika. Isi bagian kosong dengan angka
atau frasa yang sesuai.
\begin{adjustwidth}{2em}{2em}
Kita menuliskan \textit{hidup} pada \rule{2cm}{0.5pt} kartu yang mewakili pasien
yang masih hidup pada akhir penelitian, dan \textit{meninggal} pada
\rule{2cm}{0.5pt} kartu yang mewakili pasien yang sudah meninggal. Kemudian,
kartu-kartu ini dikocok dan dibagi menjadi dua kelompok: satu kelompok berukuran
\rule{2cm}{0.5pt} yang mewakili perlakuan, dan kelompok lain berukuran
\rule{2cm}{0.5pt} yang mewakili kontrol. Kita menghitung selisih antara proporsi
kartu \textit{meninggal} dalam kelompok perlakuan dan kelompok kontrol
(perlakuan - kontrol), lalu mencatat nilainya. Prosedur ini diulangi 100 kali
untuk membangun distribusi yang berpusat di \rule{2cm}{0.5pt}. Terakhir, kita
menghitung proporsi simulasi dengan selisih proporsi hasil simulasi yang
\rule{2cm}{0.5pt}. Jika proporsi ini kecil, kita menyimpulkan bahwa hasil
seperti itu kecil kemungkinannya untuk teramati secara kebetulan; hipotesis nol
harus ditolak dan hipotesis alternatif dipilih.
\end{adjustwidth}
\item Apa yang ditunjukkan oleh hasil simulasi berikut mengenai keefektifan
program transplantasi?
\end{subparts}
\end{parts}
\begin{center}
\Figures[Ditampilkan diagram titik bertumpuk yang tampaknya memuat sekitar 100 titik untuk variabel "Selisih proporsi hasil simulasi", yang membentang dari -0.25 hingga 0.25. Terdapat 11 tumpukan titik, dengan perkiraan lokasi dan jumlah sebagai berikut: 2 titik pada -0.23, 1 titik pada -0.19, 8 titik pada -0.14, 15 titik pada -0.10, 18 titik pada -0.05, 20 titik pada -0.01, 12 titik pada 0.04, 10 titik pada 0.08, 6 titik pada 0.12, 4 titik pada 0.17, dan 3 titik pada 0.21.]{0.6}{eoce/randomization_heart_transplants}{randomization_heart_transplants_rando}
\end{center}
}{}
